\section{结论}
%

我们基于对相应系数函数中阶乘发散（renormalon）的研究，在气泡链近似下分析了 qPDF 与 pPDF 的幂次压低（高扭度）贡献。阶乘渐近行为意味着级数之和仅在幂次精度内有定义，因此 QCD 微扰论必须补充非微扰的幂次压低贡献才能给出无歧义的预言。我们的结果应被视作高扭度修正的一个``最小模型''：它包含了理论自洽性所必需的效应，但可能遗漏其他非微扰修正，例如受对称性保护从而不``显现''于微扰展开中的那些。我们的主要结论如下：
\begin{itemize}
\item 位置空间 PDF（qITD）在大 Ioffe 时间处的幂修正是平缓的。一般而言，``横向''投影的幂修正远大于``纵向''投影。
\item qPDF 的幂修正具有普适行为
\begin{align}
  \cQ (x,p) &=  q(x) \biggl\{1 + \mathcal{O}\left( \frac{\Lambda^2}{p^2}\cdot\frac{1}{{x^2(1-x)}}\right) \biggr\}.
\end{align}
注意相应的靶质量修正 \eqref{target-parallel} 与 \eqref{target-perp} 并不表现出 $1/x^2$ 增强。这一行为与 DIS 结构函数中 Nachtmann 靶质量修正 $\sim x^2 m^2/Q^2$ 在小 $x$ 处的压低是相称的。将底层 qITD 按零动量归一化到一，会明显减小 $ x \lesssim 0.5$ 区域内 qPDF 的幂修正，代价是在 $x\gtrsim 0.5-0.6$ 处强烈增强。因此这种归一化步骤不适合研究 PDF 的大-$x$ 行为，但显然能使中间 $x$ 区域的幂修正最小化。注意上述讨论基于的归一化步骤不同于文献~\cite{Chen:2018xof} 中使用的非微扰重正化，后者的幂修正行为可能与本文所示不同。
\item
pPDF 的幂修正具有普适行为
\begin{align}
\mathcal{P}(x,z) &= q(x)\Big\{1 + \mathcal{O}\big( z^2\!\Lambda^2{(1-x)}\big)\Big\}.
\end{align}
$x\to 1$ 处的压低会被零动量归一化因子消除。但若按同一算符的真空矩阵元归一化，则该压低可以保持。我们由此得出结论：对于大-$x$ 区域的研究，pPDF 可以成为 qPDF 的一个有趣替代方案。
\end{itemize}
作为本研究的副产品，我们在大-$n_f$ 近似下给出了系数函数的全阶结果，即 $\sim n_f^k\alpha_s^{k+1}$ 的项。
这些结果可用于估计未计算高阶效应及采用 BLM 型步骤设定尺度的影响。
相应分析超出了本文的范围，将在未来的出版物中给出。

